\chapter{Literature Review}
\label{chap:literature}

\section{Purpose and scope}
Explain what the review covers, its boundaries and how it supports the research
questions. A literature review should synthesise relationships and disagreements,
not become a sequence of disconnected paper summaries.

\section{Search and selection approach}
Report databases, search concepts, date ranges, inclusion criteria and the final
screening process at a level another researcher could understand. If the thesis is not
a formal systematic review, describe the approach accurately rather than adopting a
label that the method cannot support.

\section{Theme one: structured scholarly writing}
Structured source documents separate content from presentation and allow the same
material to be transformed consistently \citep{lamport1994latex}. The practical
value for a large thesis is not simply attractive typesetting: stable labels,
bibliographic automation and modular files reduce avoidable maintenance work.

\section{Theme two: transparent revision}
Revision evidence should connect a specific examiner comment to a specific substantive
change. Colour can assist navigation, but identifiers such as \texttt{E1-01} preserve
meaning when a document is printed in greyscale or read with assistive technology.

\section{Synthesis and conceptual framework}
Compare the themes, explain their relationship, and end with the precise proposition
or gap taken into the methodology. \Cref{tab:literature-synthesis} demonstrates a
compact synthesis table.

\begin{table}[htbp]
  \centering
  \caption{Example literature synthesis table}
  \label{tab:literature-synthesis}
  \newcolumntype{Y}{>{\raggedright\arraybackslash}X}
  \begin{tabularx}{\textwidth}{@{}>{\raggedright\arraybackslash}p{0.18\textwidth}YYY@{}}
    \toprule
    \textbf{Theme} & \textbf{What is established} & \textbf{What is contested or missing} & \textbf{Relevance to this thesis} \\
    \midrule
    Structured writing & Modular sources improve document maintainability. & Practical workflows vary across disciplines. & Motivates the file architecture. \\
    Revision evidence & Traceable responses improve review transparency. & Colour-only systems create accessibility risks. & Motivates IDs plus colour. \\
    \bottomrule
  \end{tabularx}
\end{table}

